\section{The problem and the goal}\label{sec-problem}

\begin{objectif}
Describe the chain from the Totem contributions to the control lines, show why an implied repo curve often looks bad
without being bad, and state what the note must deliver.
\end{objectif}

\subsection{The control chain}
For each underlying and each control date, the controller has three sources (\cref{fig-chain}).
\begin{itemize}
\item \textbf{Totem options}: for each maturity $T_k$, consensus call and put prices on a strike grid, the number of
contributors $n^{\opt}_k$ and the dispersion of their quotes.
\item \textbf{Totem dividends}: for each period $p$ (usually a calendar year), the consensus level $\ell^{\tot}_p$,
that is, the sum of the dividends going ex in the period, the number of contributors $n^{\dv}_p$ and their dispersion.
\item \textbf{The front office}: the spot, the rate curve, the repo curve $q^{\fo}$, the dividend forecast (ex-dates
and amounts) and the sensitivities of each book to the repo and to the dividends.
\end{itemize}
Put--call parity gives the Totem forwards $F^{\tot}_k$. The dividend levels are spread over the FO ex-dates. A
bootstrap then gives the repo curve $q^{\tot}$ which, with these dividends, reprices the Totem forwards exactly. This
curve has two uses: computing the control lines, that is, the Totem minus FO differences times the sensitivities
(\cref{sec-control}), and checking the quality of the curve itself (\cref{sec-detection}).

\begin{figure}[ht]
\centering
\begin{tikzpicture}[
  font=\small, >={Stealth[length=2mm]},
  box/.style={rounded corners=2pt, align=center, inner sep=4pt, minimum height=1.2cm},
  src/.style={box, draw=accent, fill=accentwash, minimum width=3.3cm},
  fo/.style={box, draw=muted, fill=greywash, minimum width=3.3cm},
  mid/.style={box, draw=rule, fill=white, minimum width=2.2cm},
  key/.style={box, draw=accentdark, line width=0.9pt, fill=white, minimum width=2.2cm},
  use/.style={box, draw=green, fill=greenwash, minimum width=3.5cm},
  lab/.style={font=\footnotesize\itshape, text=ink2, align=center}]
\node[src] (opt) at (0,2.4) {\textbf{Totem options}\\ calls and puts, $n^{\opt}_k$};
\node[src] (div) at (0,0.8) {\textbf{Totem dividends}\\ levels $\ell^{\tot}_p$, $n^{\dv}_p$};
\node[fo] (fo) at (0,-1.0) {\textbf{Front office}\\ $q^{\fo}$, dividends,\\ sensitivities $s^q$, $s^\ell$};
\node[mid] (F) at (4.6,2.4) {forwards\\ $F^{\tot}_k$};
\node[mid] (D) at (4.6,0.8) {dividends\\ $D^{\tot}_k$};
\node[key] (Q) at (7.65,1.6) {\textbf{implied repo}\\ $q^{\tot}_k$};
\node[use] (L) at (11.25,2.4) {\textbf{Control lines}\\ differences $\times$ sensitivities\\ \cref{sec-control}};
\node[use] (Z) at (11.25,0.6) {\textbf{Detection}\\ on the carry difference $\delta_k$\\ \cref{sec-detection}};
\draw[->] (opt) -- node[lab, above] {parity} (F);
\draw[->] (div) -- node[lab, above] {allocation} (D);
\draw[->] (F.east) -- (Q);
\draw[->] (D.east) -- (Q);
\node[lab, anchor=south] at ($(Q.north)+(0,0.05)$) {bootstrap};
\draw[->] (Q) -- (L);
\draw[->] (Q) -- (Z);
\draw[->, dashed, draw=muted] ($(fo.north east)+(0,-0.15)$) -- (D.south);
\node[lab, anchor=north west] at ($(D.south)+(0.1,-0.12)$) {ex-dates};
\draw[draw=muted] (fo.east) -- (13.45,-1.0) -- (13.45,2.4);
\draw[->, draw=muted] (13.45,2.4) -- (L.east);
\draw[->, draw=muted] (13.45,0.6) -- (Z.east);
\node[lab, anchor=north] at (7.65,-1.05) {$q^{\fo}$, FO dividends, sensitivities};
\end{tikzpicture}
\caption{The control chain. The two Totem services have different contributors, and their number changes with
maturity; the implied repo curve $q^{\tot}$ inherits the noise of both.}
\label{fig-chain}
\end{figure}

\subsection{Why the curve looks bad}
Two reasons, neither of them an error, are enough to make a clean curve look irregular.

\paragraph{Two panels, two averages.} The Totem forward is the average of the option panel's forwards, and each
contributor's forward contains its own view of the dividends. The controller takes out the dividends of another panel,
the dividend service's. If the two panels do not share the same average view of the dividends, the difference ends
up in the implied repo. This is not a Totem error: it is the mechanical consequence of two averages taken over
different contributors. Its size depends on the contributor counts, which fall with maturity, and on the dispersion of
the dividend views, which grows with it (\cref{sec-gap}).

\paragraph{The $1/T$ magnifier.} The repo is a rate: an error $\varepsilon$ on the log of the forward becomes an error
$\varepsilon/T$ on the repo. An error of \pct{0.1} on the forward is worth \num{120}\,bp of repo at one month and
\num{2}\,bp at five years. The short end therefore always looks noisy, while a large error at the long end goes
unnoticed.

A curve that looks irregular is therefore not necessarily bad, and a curve that looks smooth is not necessarily good.
The rest of the note works on the \emph{carry difference} $\delta_k=(q^{\tot}_k-q^{\fo}_k)\,T_k$, where errors add up
without being magnified, and compares it with a noise band computed from the statistics that Totem publishes.

\subsection{The running example}\label{sec-running}
\begin{filrouge}
A synthetic European stock (\cref{sec-benchmark}): spot $S_0=100$, an annual dividend of about \vRunDiv{} going ex in
early May (the first on \vRunExDate), an FO repo going from \vRunQfoShort\,bp at the short end to \vRunQfoLong\,bp at
five years. Totem has \vRunNoptShort{} option contributors on the short maturities and \vRunNoptLong{} at five years;
the dividend service has \vRunNdivShort{} on 2027 and \vRunNdivLong{} on 2031. The eleven Totem maturities run from
one month to five years. \Cref{fig-clean} shows the clean market; the same name is taken up again, with the same
contributor noise, under eight anomalies in \cref{sec-signatures} (\cref{fig-gallery}) and run through the detector in
\cref{sec-back}.
\end{filrouge}

\begin{figure}[ht]
\centering\includegraphics[width=\linewidth]{figures/\lang/clean.pdf}
\caption{Running example, clean market. (a) FO repo and Totem-implied repo, with the band of $\pm2$ standard deviations
of consensus noise around the FO repo. (b) The same difference seen in carry: each point is
$\delta_k=(q^{\tot}_k-q^{\fo}_k)\,T_k$, in basis points of spot.}
\label{fig-clean}
\end{figure}

In repo space, the Totem curve looks badly shaped: it loses \vRunDrop\,bp between six and nine months and ends
\vRunGapLong\,bp below the FO at five years. In carry space, every point stays inside the band. The drop at nine months
corresponds to a gap of \vRunDivGap{} on the May 2027 dividend between the option panel's view and the dividend panel's,
that is, \vRunDivGapPct{} of the dividend and \vRunZNine{} standard deviations; at five years, the carry difference is
\vRunDeltaFive\,bp of spot for a standard deviation of \vRunSdFive\,bp. The detector of \cref{sec-detection} flags
nothing. This curve is clean: it is what an implied repo from two different panels normally looks like.

\subsection{What the note must deliver}
The three goals of the first page translate into four deliverables.
\begin{enumerate}
\item \textbf{A diagnosis} (G1, \cref{sec-diagnostic}): a decomposition of the implied repo that separates the
expected noise, whose size can be computed, from the mechanisms that cause real anomalies, each with a recognisable
shape.
\item \textbf{A control rule} (G2, \cref{sec-control}): a decomposition of the control lines that does not depend on
the dividend contributors where possible, and a measure, computed from the sensitivities, of what it still leaves
depending on them.
\item \textbf{A detector} (G3, \cref{sec-detection}) with a known false alarm rate, which names the mechanism at work,
gives its size in natural units (spot off by so much, May 2027 dividend too low by so much), repairs the curve, and
tells a Totem anomaly from an FO mismark.
\item \textbf{A comparison} (\cref{sec-benchmark}) with the starting ideas: B-spline coefficients by region and
category clustered by a Gaussian mixture, small neural networks specialised by anomaly type.
\end{enumerate}

\begin{retenir}
\begin{itemize}[leftmargin=12pt]
\item The implied repo curve is a residual: one panel's forwards minus another panel's dividends. It inherits the noise
of both.
\item The eye misleads: the repo magnifies short-end errors by $1/T$ and flattens long-end ones. Judge the curve on the
carry difference $\delta_k$, against a noise band.
\item In the running example, a curve that loses \vRunDrop\,bp in three months and ends \vRunGapLong\,bp below the FO
is perfectly clean.
\end{itemize}
\end{retenir}
